% =====================================================================
\tmcapitulo{tmNavy}{Las herramientas de c\'odigo que usan los ejemplos del libro, en una sola gu\'ia.}{x}
\chapter{Python para este libro}
\markboth{Ap\'endice D: Python}{}
\tmgancho{No necesitas instalar nada: abre Google Colab, crea un cuaderno nuevo y pega los bloques de c\'odigo del libro.}

\section*{Preparar el entorno}
\begin{tmpasos}
\item Entra a \texttt{colab.research.google.com} con tu cuenta de Google y elige \textbf{Cuaderno nuevo}.
\item En una celda escribe \texttt{import numpy, sympy, scipy, matplotlib} y ejec\'utala. Si no da error, todo est\'a listo.
\item Copia un bloque del libro en una celda y presiona \textbf{Shift + Enter}. Compara la salida con el comentario del bloque.
\end{tmpasos}

\begin{tmtip}Si trabajas en tu computadora, instala todo con \texttt{pip install numpy sympy scipy matplotlib}. En Windows usa el s\'imbolo del sistema o PowerShell.\end{tmtip}

\section*{Qu\'e librer\'ia usar en cada cap\'itulo}
\begin{center}
\begin{tabular}{|c|p{3.6cm}|p{7.6cm}|}
\hline
\rowcolor{tmAzulClaro}\textbf{Cap.} & \textbf{Librer\'ia} & \textbf{Funciones \'utiles}\\ \hline
1 & Python puro, NumPy & \texttt{set}, operadores \texttt{|}, \texttt{\&}, \texttt{-}, \texttt{\^{}}; \texttt{np.array}\\ \hline
2 & \texttt{math}, \texttt{sympy} & \texttt{comb}, \texttt{perm}, \texttt{factorial}, \texttt{gcd}, \texttt{pow(a, b, m)}, \texttt{factorint}, \texttt{isprime}\\ \hline
3 & NumPy, SymPy & \texttt{np.linalg.solve}, \texttt{det}, \texttt{inv}; \texttt{Matrix(...).rref()}\\ \hline
4 & SymPy, Matplotlib & \texttt{limit}, \texttt{diff}, \texttt{integrate}, \texttt{solve}; \texttt{plt.plot}\\ \hline
5 & \texttt{statistics}, SciPy & \texttt{mean}, \texttt{median}, \texttt{stdev}; \texttt{stats.binom}, \texttt{poisson}, \texttt{norm}; \texttt{np.polyfit}\\ \hline
\end{tabular}
\end{center}

\section*{Tres programas completos}
\subsection{Camino m\'as corto con Dijkstra (cap\'itulo 2)}
\begin{tmcodigo}
import heapq

G = {"A": {"B": 7, "E": 1},
     "B": {"A": 7, "C": 3, "E": 8},
     "C": {"B": 3, "D": 6, "E": 2},
     "D": {"C": 6, "E": 7},
     "E": {"A": 1, "B": 8, "C": 2, "D": 7}}

def dijkstra(G, origen):
    dist = {v: float("inf") for v in G}
    dist[origen] = 0
    cola = [(0, origen)]
    while cola:
        d, u = heapq.heappop(cola)
        if d > dist[u]:
            continue
        for v, peso in G[u].items():
            if d + peso < dist[v]:
                dist[v] = d + peso
                heapq.heappush(cola, (dist[v], v))
    return dist

print(dijkstra(G, "A"))   # {'A': 0, 'B': 6, 'C': 3, 'D': 8, 'E': 1}
\end{tmcodigo}

\subsection{Simulaci\'on de disponibilidad: m\'etodo Monte Carlo (cap\'itulo 5)}
En lugar de calcular $1-0{,}1^3$, \textbf{simulamos} 200\,000 d\'ias y contamos en cu\'antos sigue activo al menos un servidor:

\begin{tmcodigo}
import random

random.seed(1)
N = 200_000
activos = 0
for _ in range(N):
    if any(random.random() > 0.1 for _ in range(3)):   # 3 replicas
        activos += 1

print(activos / N)    # ~ 0.999 (exacto: 0.999)
\end{tmcodigo}

\subsection{Resumen estad\'istico de latencias (cap\'itulo 5)}
\begin{tmcodigo}
import statistics as st
import numpy as np
from scipy import stats

lat = [10, 12, 12, 14, 15, 16, 18, 43]
print(st.mean(lat), st.median(lat), st.stdev(lat))   # 17.5 14.5 10.61
print(np.percentile(lat, 95))                        # p95
print(stats.norm.cdf(240, loc=200, scale=20))        # 0.9772
print(stats.poisson.pmf(0, mu=3))                    # 0.0498
\end{tmcodigo}

\begin{tmcuidado}Una simulaci\'on da un resultado \textbf{aproximado} que cambia con la semilla; la f\'ormula da el valor exacto. Usa la simulaci\'on para verificar tu c\'alculo, no para sustituirlo.\end{tmcuidado}

\begin{tmejercicios}[Practica con c\'odigo]
\begin{tmlista}
\item \tmbasico\ Verifica en Python que $P(10, 3)=720$ usando \texttt{math.perm}.
\item \tmmedio\ Escribe una funci\'on que calcule el m\'aximo com\'un divisor con el algoritmo de Euclides y comp\'aralo con \texttt{math.gcd}.
\item \tmmedio\ Cambia el grafo del Dijkstra para que el peso $E$--$D$ sea 1. \textquestiondown Cambia el camino $A\to D$?
\item \tmdificil\ Simula 200\,000 d\'ias con 2 r\'eplicas en vez de 3 y compara con el valor exacto $1-0{,}1^2$.
\end{tmlista}
\end{tmejercicios}

% =====================================================================
\tmcapitulo{tmNavy}{Los errores que m\'as se repiten, reunidos para repasar antes de un examen.}{x}
\chapter{Errores comunes}
\markboth{Ap\'endice E: Errores comunes}{}
\tmgancho{Los errores no son al azar: casi siempre son los mismos. Con esta lista puedes revisar tu propio trabajo en un minuto.}

\begin{center}
\renewcommand{\arraystretch}{1.35}
\begin{tabular}{|c|p{5.4cm}|p{6.2cm}|}
\hline
\rowcolor{tmCoral}\textcolor{white}{\textbf{Cap.}} & \textcolor{white}{\textbf{Error}} & \textcolor{white}{\textbf{C\'omo evitarlo}}\\ \hline
1 & Confundir $a\in A$ con $\{a\}\subseteq A$ & Pertenencia es de elemento a conjunto; inclusi\'on es entre conjuntos\\ \hline
1 & Pensar que $p\to q$ es falsa si $p$ es falsa & Solo es falsa con $p$ verdadera y $q$ falsa\\ \hline
1 & Creer que la rec\'iproca equivale a la implicaci\'on & Solo la contrapositiva es equivalente\\ \hline
1 & Usar Big-O como igualdad exacta & Es una cota superior; busca la m\'as ajustada\\ \hline
1 & Usar un $\alpha$ muy grande en el descenso del gradiente & Si diverge, reduce $\alpha$\\ \hline
2 & Multiplicar etapas que no son independientes & Cuenta solo combinaciones v\'alidas\\ \hline
2 & Usar $P(n, r)$ cuando hay repeticiones & Con repetici\'on son $n^r$\\ \hline
2 & Llamar primo al 1 & El 1 tiene un solo divisor positivo\\ \hline
2 & Usar Dijkstra con pesos negativos & Exige pesos no negativos\\ \hline
3 & No comprobar tras elevar al cuadrado & Verifica siempre: pueden aparecer soluciones extra\~nas\\ \hline
3 & Cambiar el signo de solo un t\'ermino al quitar par\'entesis & Cambian todos los t\'erminos de adentro\\ \hline
3 & Suponer $AB=BA$ & El producto de matrices no es conmutativo\\ \hline
3 & Usar una clave de Hill con $\det$ no coprimo con 26 & Verifica $\operatorname{MCD}(\det K, 26)=1$\\ \hline
4 & Olvidar la cadena al derivar & Multiplica por la derivada de lo de adentro\\ \hline
4 & Dar por m\'inimo todo punto con $f'=0$ & Clasifica con $f''$ (o con $D$ en dos variables)\\ \hline
4 & Omitir la constante $+C$ & En la integral indefinida siempre va\\ \hline
5 & Confundir $P(A\mid B)$ con $P(B\mid A)$ & Usa Bayes\\ \hline
5 & Esperar que $E[X]$ sea un resultado posible & Es un promedio a largo plazo\\ \hline
5 & Dividir entre $n$ la varianza muestral & Se divide entre $n-1$\\ \hline
5 & Concluir causalidad de una correlaci\'on & Correlaci\'on no implica causalidad\\ \hline
\end{tabular}
\end{center}

\begin{tmtip}Cuando termines un problema, rev\'isalo con esta pregunta: \textquestiondown qu\'e error de esta tabla podr\'ia haber cometido? Si no encuentras ninguno, probablemente est\'a bien.\end{tmtip}

\begin{tmquiz}[Mini repaso: \textquestiondown qu\'e error es?]
\tmpreg{Alguien dice que $\sqrt{x+7}=x+1$ tiene soluciones $x=2$ y $x=-3$.}{No comprob\'o la ra\'iz extra\~na}{Olvid\'o la cadena}{Us\'o $P(n, r)$}{Divid\'o entre $n$}
\tmpreg{Alguien calcula $(2x+1)^3$ derivado como $3(2x+1)^2$.}{Olvid\'o la constante $C$}{Olvid\'o la regla de la cadena}{Us\'o Bayes mal}{No clasific\'o con $f''$}
\tmpreg{Alguien dice: ``si el correo es spam contiene la palabra con 60\%, entonces un correo con la palabra es spam con 60\%''.}{Confundi\'o $P(A\mid B)$ con $P(B\mid A)$}{Us\'o conteo}{Olvid\'o $+C$}{Us\'o Dijkstra}
\end{tmquiz}
\textbf{Respuestas:} 1 a, 2 b, 3 a.

